\documentclass[10pt,a4paper]{amsart}
\usepackage{amsmath,amssymb,amsthm,graphicx}
\usepackage{verbatim}
\newenvironment{reference}{\comment}{\endcomment}
%\newenvironment{reference}{}{}
\newenvironment{slogan}{\comment}{\endcomment}
\newenvironment{history}{\comment}{\endcomment}

% For commutative diagrams we use Xy-pic
\usepackage[all]{xy}

% We use 2cell for 2-commutative diagrams.
\xyoption{2cell}
% Theorem environments.
%
\theoremstyle{plain}
\newtheorem{theorem}[subsection]{Theorem}
\newtheorem{proposition}[subsection]{Proposition}
\newtheorem{lemma}[subsection]{Lemma}

\theoremstyle{definition}
\newtheorem{definition}[subsection]{Definition}
\newtheorem{example}[subsection]{Example}
\newtheorem{exercise}[subsection]{Exercise}
\newtheorem{situation}[subsection]{Situation}

\theoremstyle{remark}
\newtheorem{remark}[subsection]{Remark}
\newtheorem{remarks}[subsection]{Remarks}

\numberwithin{equation}{subsection}

% Macros
%
\def\lim{\mathop{\mathrm{lim}}\nolimits}
\def\colim{\mathop{\mathrm{colim}}\nolimits}
\def\Spec{\mathop{\mathrm{Spec}}}
\def\Hom{\mathop{\mathrm{Hom}}\nolimits}
\def\Ext{\mathop{\mathrm{Ext}}\nolimits}
\def\SheafHom{\mathop{\mathcal{H}\!\mathit{om}}\nolimits}
\def\SheafExt{\mathop{\mathcal{E}\!\mathit{xt}}\nolimits}
\def\Sch{\mathit{Sch}}
\def\Mor{\mathop{\mathrm{Mor}}\nolimits}
\def\Ob{\mathop{\mathrm{Ob}}\nolimits}
\def\Sh{\mathop{\mathit{Sh}}\nolimits}
\def\NL{\mathop{N\!L}\nolimits}
\def\CH{\mathop{\mathrm{CH}}\nolimits}
\def\proetale{{pro\text{-}\acute{e}tale}}
\def\etale{{\acute{e}tale}}
\def\QCoh{\mathit{QCoh}}
\def\Ker{\mathop{\mathrm{Ker}}}
\def\Im{\mathop{\mathrm{Im}}}
\def\Coker{\mathop{\mathrm{Coker}}}
\def\Coim{\mathop{\mathrm{Coim}}}

% Boxtimes
%
\DeclareMathSymbol{\boxtimes}{\mathbin}{AMSa}{"02}

%
% Macros for moduli stacks/spaces
%
\def\QCohstack{\mathcal{QC}\!\mathit{oh}}
\def\Cohstack{\mathcal{C}\!\mathit{oh}}
\def\Spacesstack{\mathcal{S}\!\mathit{paces}}
\def\Quotfunctor{\mathrm{Quot}}
\def\Hilbfunctor{\mathrm{Hilb}}
\def\Curvesstack{\mathcal{C}\!\mathit{urves}}
\def\Polarizedstack{\mathcal{P}\!\mathit{olarized}}
\def\Complexesstack{\mathcal{C}\!\mathit{omplexes}}
% \Pic is the operator that assigns to X its picard group, usage \Pic(X)
% \Picardstack_{X/B} denotes the Picard stack of X over B
% \Picardfunctor_{X/B} denotes the Picard functor of X over B
\def\Pic{\mathop{\mathrm{Pic}}\nolimits}
\def\Picardstack{\mathcal{P}\!\mathit{ic}}
\def\Picardfunctor{\mathrm{Pic}}
\def\Deformationcategory{\mathcal{D}\!\mathit{ef}}
\usepackage{cleveref}
\crefname{originallemma}{original lemma}{original lemmas}
\Crefname{originallemma}{Original lemma}{Original lemmas}
\makeatletter
\newlabel{lemma-general-pushforward}{{076Z}{source}{}{sourcepointer-1}{}}
\newlabel{lemma-general-pushforward@cref}{{[originallemma][0][]076Z}{[1][0][]source}{}{}{}}
\makeatother
\crefname{originalreference}{original reference}{original references}
\Crefname{originalreference}{Original reference}{Original references}
\makeatletter
\newlabel{stacks-sheaves-lemma-pullback-lqc}{{06WL}{source}{}{sourcepointer-2}{}}
\newlabel{stacks-sheaves-lemma-pullback-lqc@cref}{{[originalreference][0][]06WL}{[1][0][]source}{}{}{}}
\makeatother
\crefname{originalreference}{original reference}{original references}
\Crefname{originalreference}{Original reference}{Original references}
\makeatletter
\newlabel{stacks-sheaves-lemma-lqc-colimits}{{06WM}{source}{}{sourcepointer-3}{}}
\newlabel{stacks-sheaves-lemma-lqc-colimits@cref}{{[originalreference][0][]06WM}{[1][0][]source}{}{}{}}
\makeatother
\crefname{originallemma}{original lemma}{original lemmas}
\Crefname{originallemma}{Original lemma}{Original lemmas}
\makeatletter
\newlabel{lemma-check-lqc-on-etale-covering}{{075Y}{source}{}{sourcepointer-4}{}}
\newlabel{lemma-check-lqc-on-etale-covering@cref}{{[originallemma][0][]075Y}{[1][0][]source}{}{}{}}
\makeatother
\crefname{originallemma}{original lemma}{original lemmas}
\Crefname{originallemma}{Original lemma}{Original lemmas}
\makeatletter
\newlabel{lemma-check-flat-comparison-on-etale-covering}{{0764}{source}{}{sourcepointer-5}{}}
\newlabel{lemma-check-flat-comparison-on-etale-covering@cref}{{[originallemma][0][]0764}{[1][0][]source}{}{}{}}
\makeatother
\crefname{originalreference}{original reference}{original references}
\Crefname{originalreference}{Original reference}{Original references}
\makeatletter
\newlabel{stacks-sheaves-lemma-compare-representable-morphism-cohomology}{{075P}{source}{}{sourcepointer-6}{}}
\newlabel{stacks-sheaves-lemma-compare-representable-morphism-cohomology@cref}{{[originalreference][0][]075P}{[1][0][]source}{}{}{}}
\makeatother
\crefname{originalreference}{original reference}{original references}
\Crefname{originalreference}{Original reference}{Original references}
\makeatletter
\newlabel{spaces-cohomology-lemma-flat-base-change-cohomology}{{073K}{source}{}{sourcepointer-7}{}}
\newlabel{spaces-cohomology-lemma-flat-base-change-cohomology@cref}{{[originalreference][0][]073K}{[1][0][]source}{}{}{}}
\makeatother
\crefname{originallemma}{original lemma}{original lemmas}
\Crefname{originallemma}{Original lemma}{Original lemmas}
\makeatletter
\newlabel{lemma-pushforward-locally-quasi-coherent}{{075Z}{source}{}{sourcepointer-8}{}}
\newlabel{lemma-pushforward-locally-quasi-coherent@cref}{{[originallemma][0][]075Z}{[1][0][]source}{}{}{}}
\makeatother
\title{Higher direct images of lqc fbc modules have the flat base change property}
\author{The Stacks Project Authors}
\date{stacks-project:stacks-cohomology.tex:local-snapshot-20261012; GNU FDL 1.2 or later}
\begin{document}
\maketitle
\noindent\textit{Editorial note.} All delivered mathematical text is the author's own, copied verbatim from the pinned Stacks Project chapter file stacks-cohomology.tex; no statement or proof step was rewritten, reordered or removed.\par\medskip
\noindent\textit{Editorial note.} The counted span cites the lemmas recorded in assembly.external\_references with canonical tags as navigation pointers.\par\medskip
\noindent\textit{Allowed standard prerequisites (editorial).}\begin{itemize}
\item General pushforward criterion. Cohomology of Algebraic Stacks, Lemma stacks-cohomology-lemma-general-pushforward. Supplies the four conditions whose verification yields the theorem.
\item Comparison of cohomology for representable morphisms. Sheaves on Stacks, Lemma stacks-sheaves-lemma-compare-representable-morphism-cohomology. Identifies R\textasciicircum{}i f\_* F with the pushforward of the restriction to the scheme cover.
\item Flat base change for cohomology of algebraic spaces. Cohomology of Spaces, Lemma spaces-cohomology-lemma-flat-base-change-cohomology. Gives the isomorphism of the comparison map in the scheme case.
\end{itemize}\medskip
\ifdefined\hypertarget\hypertarget{sourcepointer-1}{}\fi
\noindent\textit{Source reference.} Original lemma 076Z: sources/stacks-chow/tags.txt, source lines 9184--9184. Cross-chapter standard prerequisite cited by the author inside the counted span; it is used as a pointer and its own proof is not part of this item.\par\medskip
\ifdefined\hypertarget\hypertarget{sourcepointer-2}{}\fi
\noindent\textit{Source reference.} Original reference 06WL: sources/stacks-chow/tags.txt, source lines 8820--8820. Cross-chapter standard prerequisite cited by the author inside the counted span; it is used as a pointer and its own proof is not part of this item.\par\medskip
\ifdefined\hypertarget\hypertarget{sourcepointer-3}{}\fi
\noindent\textit{Source reference.} Original reference 06WM: sources/stacks-chow/tags.txt, source lines 8821--8821. Cross-chapter standard prerequisite cited by the author inside the counted span; it is used as a pointer and its own proof is not part of this item.\par\medskip
\ifdefined\hypertarget\hypertarget{sourcepointer-4}{}\fi
\noindent\textit{Source reference.} Original lemma 075Y: sources/stacks-chow/tags.txt, source lines 9147--9147. Cross-chapter standard prerequisite cited by the author inside the counted span; it is used as a pointer and its own proof is not part of this item.\par\medskip
\ifdefined\hypertarget\hypertarget{sourcepointer-5}{}\fi
\noindent\textit{Source reference.} Original lemma 0764: sources/stacks-chow/tags.txt, source lines 9154--9154. Cross-chapter standard prerequisite cited by the author inside the counted span; it is used as a pointer and its own proof is not part of this item.\par\medskip
\ifdefined\hypertarget\hypertarget{sourcepointer-6}{}\fi
\noindent\textit{Source reference.} Original reference 075P: sources/stacks-chow/tags.txt, source lines 9138--9138. Cross-chapter standard prerequisite cited by the author inside the counted span; it is used as a pointer and its own proof is not part of this item.\par\medskip
\ifdefined\hypertarget\hypertarget{sourcepointer-7}{}\fi
\noindent\textit{Source reference.} Original reference 073K: sources/stacks-chow/tags.txt, source lines 9064--9064. Cross-chapter standard prerequisite cited by the author inside the counted span; it is used as a pointer and its own proof is not part of this item.\par\medskip
\ifdefined\hypertarget\hypertarget{sourcepointer-8}{}\fi
\noindent\textit{Source reference.} Original lemma 075Z: sources/stacks-chow/tags.txt, source lines 9148--9148. Cross-chapter standard prerequisite cited by the author inside the counted span; it is used as a pointer and its own proof is not part of this item.\par\medskip
\begin{lemma}
\label{lemma-flat-comparison}
Let $f : \mathcal{X} \to \mathcal{Y}$ be a quasi-compact and
quasi-separated morphism of algebraic stacks. Let 
$\mathcal{F}$ be an object of
$\textit{Mod}(\mathcal{X}_\etale, \mathcal{O}_\mathcal{X})$
which is locally quasi-coherent and has the flat base change property.
Then each $R^if_*\mathcal{F}$ (computed in the \'etale topology)
has the flat base change property.
\end{lemma}

\begin{proof}
We will use
Lemma \ref{lemma-general-pushforward}
to prove this. For every algebraic stack $\mathcal{X}$ let
$\textit{LQCoh}^{fbc}(\mathcal{O}_\mathcal{X})$ denote the full subcategory of
$\textit{Mod}(\mathcal{X}_\etale, \mathcal{O}_\mathcal{X})$
consisting of locally quasi-coherent sheaves with the flat base
change property. Once we verify conditions (1) -- (4) of
Lemma \ref{lemma-general-pushforward}
the lemma will follow. Properties (1), (2), and (3) follow from
Sheaves on Stacks, Lemmas \ref{stacks-sheaves-lemma-pullback-lqc} and
\ref{stacks-sheaves-lemma-lqc-colimits}
and
Lemmas \ref{lemma-check-lqc-on-etale-covering} and
\ref{lemma-check-flat-comparison-on-etale-covering}.
Thus it suffices to show part (4).

\medskip\noindent
Suppose $f : \mathcal{X} \to \mathcal{Y}$ is a morphism of algebraic stacks
such that $\mathcal{X}$ and $\mathcal{Y}$ are representable by affine
schemes $X$ and $Y$. In this case, suppose that
$\psi : y \to y'$ is a morphism of $\mathcal{Y}$ lying over
a flat morphism $b : V \to V'$ of schemes. For clarity denote
$\mathcal{V} = (\Sch/V)_{fppf}$ and $\mathcal{V}' = (\Sch/V')_{fppf}$
the corresponding algebraic stacks. Consider the diagram
of algebraic stacks
$$
\xymatrix{
\mathcal{Z} \ar[d]_{f''} \ar[r]_a &
\mathcal{Z}' \ar[r]_{x'} \ar[d]_{f'} & \mathcal{X} \ar[d]^f \\
\mathcal{V} \ar[r]^b & \mathcal{V}' \ar[r]^{y'} & \mathcal{Y}
}
$$
with both squares cartesian. As $f$ is representable by schemes
(and quasi-compact and separated -- even affine) we see that $\mathcal{Z}$ and
$\mathcal{Z}'$ are representable by schemes $Z$ and $Z'$ and in
fact $Z = V \times_{V'} Z'$. Since $\mathcal{F}$ has the flat
base change property we see that
$$
a_{small}^*\big(\mathcal{F}|_{Z'_\etale}\big)
\longrightarrow
\mathcal{F}|_{Z_\etale}
$$
is an isomorphism. Moreover,
$$
R^if_*\mathcal{F}|_{V'_\etale} =
R^i(f')_{small, *}\big(\mathcal{F}|_{Z'_\etale}\big)
$$
and
$$
R^if_*\mathcal{F}|_{V_\etale} =
R^i(f'')_{small, *}\big(\mathcal{F}|_{Z_\etale}\big)
$$
by
Sheaves on Stacks, Lemma
\ref{stacks-sheaves-lemma-compare-representable-morphism-cohomology}.
Hence we see that the comparison map
$$
c_\psi :
b_{small}^*(R^if_*\mathcal{F}|_{V'_\etale})
\longrightarrow
R^if_*\mathcal{F}|_{V_\etale}
$$
is an isomorphism by
Cohomology of Spaces, Lemma
\ref{spaces-cohomology-lemma-flat-base-change-cohomology}.
Thus $R^if_*\mathcal{F}$ has the flat base change property.
Since $R^if_*\mathcal{F}$ is locally quasi-coherent by
Lemma \ref{lemma-pushforward-locally-quasi-coherent}
we win.
\end{proof}

\end{document}
